\chapter{CONCLUSIONES Y RECOMENDACIONES}
\label{ch:conclusiones}

Este capítulo debe cerrar la cadena entre problema, objetivos, resultados y contribución. Debido a que el plan no contiene evidencia final de construcción y evaluación, no se formulan conclusiones anticipadas. La estructura queda preparada para incorporar afirmaciones sustentadas después de completar el Capítulo~\ref{ch:propuesta-desarrollo}.

\section{CONCLUSIONES}

\pendiente{Redactar las conclusiones con base exclusiva en resultados verificados. Cada conclusión debe responder un objetivo, incluir la evidencia cuantitativa necesaria y declarar el dominio de validez.}

La versión final debe mantener el siguiente orden lógico:

\begin{enumerate}
  \item conclusión sobre el problema general y la utilidad demostrada del artefacto;
  \item conclusión sobre la especificación trazable y los requisitos del OE1;
  \item conclusión sobre la verificación, exactitud, consistencia física y robustez del OE2;
  \item conclusión sobre el efecto del contraste, patrón, extensión y proximidad en el OE3;
  \item conclusión sobre las discrepancias de ampacidad identificadas en el OE4;
  \item conclusión metodológica sobre los principios de diseño, la reproducibilidad y los límites del artefacto.
\end{enumerate}

Las conclusiones no deben repetir tablas ni presentar información nueva. Tampoco deben convertir una verificación computacional en validación de campo. Cuando un objetivo no se alcance, la conclusión debe explicar el resultado y su efecto sobre el alcance de la tesis.

\section{RECOMENDACIONES}

\pendiente{Formular recomendaciones después de cerrar las conclusiones. Cada recomendación debe derivarse de una limitación, un resultado o una condición observada durante la construcción y evaluación.}

Las recomendaciones podrán dirigirse a cuatro ámbitos, siempre que exista evidencia:

\begin{enumerate}
  \item uso del artefacto dentro del dominio verificado;
  \item mejoras de arquitectura, muestreo, pérdidas o tratamiento de interfaces;
  \item ampliación a régimen transitorio, humedad acoplada o validación experimental;
  \item prácticas de trazabilidad y verificación para futuros artefactos térmicos.
\end{enumerate}

No deben recomendarse decisiones sobre instalaciones reales, aumentos de capacidad ni sustitución de métodos normativos sin una validación independiente que supere el alcance declarado.

